\chapter{Diseño de índices auxiliares}
\label{cap:indices-auxiliares}

\begin{intuicion}
NLF y k-hop reachability son \emph{índices auxiliares}: estructuras
derivadas del grafo que el matcher consulta para podar trabajo. La
distinción ``auxiliar'' importa: NO son la \emph{source of truth};
si se invalidan o se vuelven inconsistentes, el matcher sigue
funcionando ---sólo más lento, sin las podas---. Este capítulo
formaliza el contrato general que cualquier índice auxiliar de
GraphHunter debe cumplir: no source of truth, soundness por
sobre-aproximación, lazy build, y un layout flexible
heap/mmap-spilled. Después los próximos dos capítulos
(\ref{cap:nlf-detalle} y \ref{cap:khop-detalle}) los aplican a NLF
y k-hop específicamente.
\end{intuicion}

\section{El contrato general}

Un índice auxiliar en GraphHunter debe cumplir cuatro propiedades:

\begin{enumerate}
    \item \textbf{No es source of truth.} La información canónica
    vive en el \emph{streaming backend} (capítulo
    \ref{cap:streaming-mirror}). El índice puede invalidarse,
    perderse, descartarse ---el motor sigue funcionando. La única
    consecuencia de no tener el índice es que el matcher no puede
    usar ese prune.
    \item \textbf{Sound como sobre-aproximación.} El índice puede
    \emph{sobre-aproximar}: incluir candidatos que después no van a
    matchear. NUNCA puede \emph{sub-aproximar}: excluir candidatos
    que sí matchearían. Eso garantiza que el prune sea
    \emph{correcto sin importar el estado del índice}: si el índice
    dice ``no'', la respuesta es \emph{no}; si dice ``sí'', el
    matcher verifica.
    \item \textbf{Lazy build.} El índice se construye al primer
    consumer que lo necesita, no eagerly al cargar la sesión.
    Implementación con \texttt{OnceLock<T>}.
    \item \textbf{Heap o mmap-spilled.} La estructura física del
    índice puede vivir en heap (mutable, rápida) o mmap-spilled
    (read-only, fría) según presione la memoria. El \emph{wrapper}
    \texttt{IndexBacking<T>} encapsula esa dualidad.
\end{enumerate}

\section{Por qué ``no source of truth''}

Una pregunta natural: si el índice acelera el matcher, ¿por qué no
hacerlo source of truth y eliminar el \emph{streaming backend}?

Tres razones que justifican la separación:

\begin{enumerate}
    \item \textbf{Lossiness.} Los índices son \emph{derivaciones
    lossy} del grafo. NLF guarda counts pero no las aristas
    individuales; k-hop guarda bits pero no los caminos concretos.
    Reconstruir el grafo desde el índice es imposible; reconstruir
    el índice desde el grafo es siempre posible.
    \item \textbf{Decoupling de schema.} Si en el futuro se cambia
    el layout de NLF (más tipos, más resolución de count), el grafo
    no se toca. Si NLF fuera source of truth, cualquier cambio de
    schema obligaría una migración.
    \item \textbf{Auditoría.} El analista puede inspeccionar
    \emph{las aristas concretas} que llevaron a un match
    (capítulo~\ref{cap:hunt-result}). Sólo el grafo canónico
    permite eso; un índice agregado lo perdería.
\end{enumerate}

\section{Soundness por sobre-aproximación}

\begin{definition}[Soundness de prune]
Sea \(\mathrm{Index}\) un índice auxiliar y \(\mathrm{prune}_{H}\)
una función booleana que el matcher invoca para descartar
candidatos para la hipótesis \(H\). El prune es \emph{sound} sii:
\[
\mathrm{prune}_H(\mathrm{Index}, c) = \texttt{true} \implies c\;\text{no genera match para}\;H.
\]
\end{definition}

\noindent
Lo equivalente: el conjunto de candidatos válidos está contenido
en el conjunto que el prune \emph{no} descarta. El prune puede
ser conservador (descartar pocos), nunca agresivo en exceso.

\begin{nota}
NLF y k-hop ambos satisfacen soundness por construcción: NLF cuenta
los vecinos \emph{exactamente}; descartar cuando el count es 0
descarta sólo donde no hay vecinos. k-hop bit es 0 sii NO existe
walk \(\leq \texttt{MAX\_DEPTH}\); descartar es seguro porque ningún
walk implica ningún match.

La sobre-aproximación entra cuando los índices \emph{podrían}
sub-aproximar pero deciden no hacerlo: por ejemplo, k-hop podría
ser ``más preciso'' si guardara \emph{cuántos} caminos hay en vez
de un bit, pero entonces sería incorrecto si el contador
overflowara (a 0). Saturating + bit es la versión segura.
\end{nota}

\section{Lazy build con \texttt{OnceLock}}

El patrón:

\begin{lstlisting}[language=Rust]
pub struct GraphHunter {
    streaming: Arc<InMemoryStreamingBackend>,
    nlf: OnceLock<NlfTable>,
    khop_reach: OnceLock<KHopReach>,
    // ...
}

impl GraphHunter {
    pub fn ensure_nlf(&self) -> &NlfTable {
        self.nlf.get_or_init(|| NlfTable::build(self))
    }

    pub fn ensure_khop_reach(&self) -> &KHopReach {
        self.khop_reach.get_or_init(|| KHopReach::build(self))
    }
}
\end{lstlisting}

\noindent
Cada \texttt{ensure\_X} construye el índice en el primer call;
los subsiguientes ven el resultado vía atómica
(capítulo~\ref{cap:sistemas}). El consumer del índice (el matcher)
nunca tiene que preocuparse por ``¿está construido?''; siempre
llama \texttt{ensure} y obtiene una referencia.

\subsection{Invalidation}

Cuando una nueva arista llega \emph{post}-build, el índice puede
quedar \emph{stale}. Tres políticas posibles:

\begin{enumerate}
    \item \textbf{Invalidación}: \texttt{self.nlf = OnceLock::new()};
    el próximo \texttt{ensure\_nlf} reconstruye desde cero.
    \item \textbf{Update incremental}: aplicar el delta sin
    rebuild. Para NLF es directo (\texttt{try\_record\_edge},
    capítulo~\ref{cap:incremental-nlf}); para k-hop es difícil
    (la transitividad propaga el cambio, capítulo~\ref{cap:khop-delta-hard}).
    \item \textbf{Best-effort + fallback}: intentar update; si
    falla por estructura física (mmap-spilled), invalidar.
\end{enumerate}

\noindent
GraphHunter usa la 3 (NLF, capítulo~\ref{cap:incremental-nlf}). La
1 era el comportamiento previo a stage 2a; la 2 pura es el plan
para k-hop (capítulo~\ref{cap:delta-wcoj}).

\section{La dualidad heap / mmap-spilled}

Los índices grandes (k-hop especialmente, varios GB) pueden no
caber en heap. El wrapper \texttt{IndexBacking<T>} encapsula la
dualidad heap-resident vs mmap-spilled, el mismo patrón que el
\emph{streaming backend} aplica a las aristas
(capítulo~\ref{cap:streaming-mirror}):

\begin{lstlisting}[language=Rust]
pub enum IndexBacking<T: Pod> {
    Heap(Vec<T>),                              // mutable
    Mmap { file: File, map: Mmap, slice: &'static [T] },  // read-only
}

impl<T: Pod> IndexBacking<T> {
    pub fn get(&self, idx: usize) -> Option<&T> {
        match self {
            Self::Heap(v) => v.get(idx),
            Self::Mmap { slice, .. } => slice.get(idx),
        }
    }

    pub fn heap_slice_mut(&mut self) -> Option<&mut [T]> {
        match self {
            Self::Heap(v) => Some(v.as_mut_slice()),
            Self::Mmap { .. } => None,  // read-only
        }
    }
}
\end{lstlisting}

\noindent
Tres notas:

\begin{itemize}
    \item \texttt{T: Pod} (\emph{Plain Old Data}): garantiza que el
    layout de \texttt{T} sea binario plano, sin pointers internos.
    Permite mmap directo.
    \item \texttt{get()} es uniforme en ambas variantes: el caller
    no sabe si lee de heap o mmap.
    \item \texttt{heap\_slice\_mut()} sólo funciona en heap;
    devuelve \texttt{None} en mmap. El caller decide qué hacer
    (invalidar, abortar update, etc.).
\end{itemize}

\subsection{Cuándo se decide spillear}

\texttt{IndexBacking::build} decide al construir:

\begin{lstlisting}[language=Rust]
pub fn build_with_budget<F>(
    capacity: usize,
    budget: usize,
    populate: F,
) -> Result<Self, Error>
where
    F: FnOnce(&mut Vec<T>),
{
    let bytes_estimated = capacity * std::mem::size_of::<T>();
    let mut v: Vec<T> = Vec::with_capacity(capacity);
    populate(&mut v);
    if bytes_estimated > budget {
        let file = create_temp_file()?;
        write_to_file(&file, &v)?;
        let map = unsafe { Mmap::map(&file)? };
        Ok(Self::Mmap { file, map, slice: ... })
    } else {
        Ok(Self::Heap(v))
    }
}
\end{lstlisting}

\noindent
El budget default por índice es 1~GB. Para grafos chico-medianos
todos los índices viven en heap; para grafos grandes (k-hop > 4
GB), el k-hop spillea pero NLF (que es chico) suele caber.

\section{El lifecycle completo de un índice}

\begin{figure}[H]
\centering
\begin{tikzpicture}[
    node distance=0.6cm,
    every node/.style={font=\scriptsize, align=center},
    state/.style={draw=accent, thick, rounded corners=2mm,
                  minimum height=0.9cm, minimum width=2.4cm,
                  fill=accentlight}
]
\node[state] (sealed) {grafo sealed\\\itshape OnceLock::new()};
\node[state, right=of sealed] (build) {build (lazy)\\\itshape primer ensure\_X};
\node[state, right=of build] (heap) {heap-resident\\\itshape mutable};
\node[state, below=of heap] (mmap) {mmap-spilled\\\itshape read-only};
\node[state, below=of build] (inval) {invalidated\\\itshape OnceLock::new() de nuevo};

\draw[->, >=Latex, accent, thick] (sealed) -- (build);
\draw[->, >=Latex, accent, thick] (build) -- node[above]{\scriptsize fits} (heap);
\draw[->, >=Latex, accent, thick] (build) -- node[left]{\scriptsize spill} (mmap);
\draw[->, >=Latex, accent, thick] (heap) -- node[right]{\scriptsize add\_relation} (inval);
\draw[->, >=Latex, accent, thick] (mmap) to[bend left=20] node[below, pos=0.4]{\scriptsize add\_relation} (inval);
\draw[->, >=Latex, accent, thick] (inval) to[bend left=20] node[above, pos=0.7]{\scriptsize next ensure\_X} (build);
\end{tikzpicture}
\caption{Lifecycle de un índice auxiliar. Arranca \texttt{None}
(sealed graph). Primer \texttt{ensure} construye y elige
heap/mmap. \texttt{add\_relation} \emph{post}-build invalida (a
menos que se logre incremental, NLF caso).}
\label{fig:index-lifecycle}
\end{figure}

\section{Ejemplo concreto: cuánto cuesta cada estado}

Para un grafo de 28 GB EVTX (200 M aristas, 30 M ranked vertices):

\begin{center}
\renewcommand{\arraystretch}{1.2}
\begin{tabular}{lll}
\toprule
\textbf{Estado} & \textbf{Tiempo} & \textbf{Memoria} \\
\midrule
sealed (\texttt{None}) & instantáneo & 0 \\
build NLF & $\sim$30 s & $\sim$1.6 GB heap \\
build k-hop & $\sim$1.5 min & $\sim$5 GB (4.5 spilled, 0.5 heap) \\
ensure (post-build) & $\sim$5 ns & ya construido \\
add\_relation incremental NLF (heap) & +50 ns & sin cambio \\
add\_relation invalidation NLF (mmap) & instantáneo & libera mmap \\
\bottomrule
\end{tabular}
\end{center}

\section{Resumen}

\begin{itemize}
    \item Índice auxiliar = derivación lossy del grafo, NO source
    of truth, soundness por sobre-aproximación.
    \item Lazy build via \texttt{OnceLock}: cero costo hasta el
    primer hunt que lo consulta.
    \item Dualidad heap / mmap-spilled vía \texttt{IndexBacking<T>}.
    \item Invalidation: simple por defecto; incremental cuando se
    puede (NLF); plan para k-hop (capítulo~\ref{cap:delta-wcoj}).
\end{itemize}

\begin{ejercicios}
\begin{enumerate}
    \item \starred{} Si el motor pierde el contenido de NLF
    (digamos, por un crash del proceso), ¿se rompe la sesión? ¿Qué
    ocurre en el próximo hunt?

    \item ¿Por qué la propiedad ``no source of truth'' es
    importante para auditoría? Pensá: si NLF fuera source of truth,
    ¿qué pregunta del analista no se podría responder?

    \item \starred{} Diseñá un índice auxiliar \emph{nuevo} para
    GraphHunter: ``\texttt{first\_seen[v]}'' = timestamp de la
    primera arista incidente a \(v\). ¿Cómo lo encajás en el
    contrato? ¿Cómo se actualiza incrementalmente?

    \item ¿Cuándo conviene NO usar \texttt{OnceLock} para un
    índice auxiliar? (Pista: pensá en índices que se usen muy
    raramente.)

    \item Si tu grafo tiene 100 M vértices y NLF es
    \texttt{[u32; 9]} por vértice, ¿cuánta memoria? ¿Excede el
    budget default de 1 GB?
\end{enumerate}
\end{ejercicios}
